
\textbf{Corollaire 3.4.} \uline{Soient} @@M0@@ \uline{un corps,} @@M1@@ \uline{un} @@M2@@\uline{-schéma propre. On suppose que} @@M3@@ \uline{possède un faisceau inversible universel. (Par exemple, ce sera le cas, d'après la théorie de descente, si toute composante connexe} @@M4@@ \uline{de} @@M5@@ \uline{possède un point rationnel, par exemple si} @@M6@@ \uline{est algébriquement clos, ou plus généralement si le} p.g.c.d. \uline{des degrés des} @@M7@@\uline{-cycles de tout} @@M8@@ \uline{est} @@M9@@\uline{). Alors au-dessus de tout ouvert affine} @@M10@@ \uline{de} @@M11@@\uline{, il existe un faisceau inversible universel sur} @@M12@@.

Finalement, on démontre I sous la forme plus précise suivante~:

\textbf{Proposition 3.5.} \uline{Soient} @@M13@@ \uline{un schéma noethérien intègre,} @@M14@@ \uline{une immersion nilpotente de} @@M15@@\uline{-schémas propres. Alors il existe un ouvert non vide} @@M16@@ \uline{de} @@M17@@ \uline{tel que le morphisme}
@@M18@@
\uline{soit représentable par un morphisme affine de présentation finie~; en fait,} @@M19@@ \uline{est un schéma en groupes affines à fibres unipotentes (au sens de} SGA 3 XVII).

En effet @@M20@@ est défini par un idéal nilpotent @@M21@@. Alors @@M22@@ possède une filtration par des idéaux @@M23@@ tels que @@M24@@, où @@M25@@ est le nilradical de @@M26@@ (par exemple @@M27@@). Soit @@M28@@ le sous-schéma de @@M29@@ défini par @@M30@@. Par récurrence il suffit d'étudier les immersions @@M31@@. On peut donc supposer que @@M32@@~; c'est la condition i) de 3.1 et on finit en se ramenant au cas où les autres conditions (ii) à (v) sont vérifiées également. Grâce au théorème de la platitude générique (EGA IV 6.9.1) appliqué à @@M33@@, et aux @@M34@@ et à @@M35@@, et compte tenu des relations de Künneth EGA III 6.9.9, on peut supposer, quitte à restreindre @@M36@@, que les conditions ii), iii) et iv) sont vérifiées.

